\documentclass[11pt]{article}
\usepackage[a4paper,margin=1in]{geometry}
\usepackage{fontspec}
\usepackage[boldfont=false]{xeCJK}
\setCJKmainfont{FandolSong-Regular.otf}
\usepackage{amsmath}
\usepackage{amssymb}
\usepackage{array}
\usepackage{booktabs}
\usepackage{graphicx}
\usepackage{adjustbox}
\usepackage{enumitem}
\setlist[itemize]{topsep=2pt,partopsep=0pt,parsep=0pt,itemsep=2pt,leftmargin=1.4em}
\setlist[enumerate]{topsep=2pt,partopsep=0pt,parsep=0pt,itemsep=2pt,leftmargin=1.6em}
\usepackage{parskip}
\usepackage{fvextra}
\fvset{breaklines=true,breakanywhere=true,fontsize=\small}
\setlength{\parindent}{0pt}

\begin{document}

\begin{center}
  {\LARGE\bfseries Superposition}\\[0.35em]
  {\large A Level Physics}
\end{center}
\vspace{0.6em}

\section*{Principle of superposition}

When two or more \textbf{waves} 波 overlap at a point, the \textbf{displacement} 位移 there is the \textbf{vector sum} 矢量和 of the displacements each wave would make on its own. This is the \textbf{principle of superposition} 叠加.

The waves pass through each other and come out unchanged. Superposition is the base of everything in this topic.

If two waves of \textbf{amplitude} 振幅 $A_{1}$ and $A_{2}$ meet:

\begin{itemize}
\item \textbf{in phase} 同相 (\textbf{crest} 波峰 meets crest): the amplitude is $A_{1} + A_{2}$ (\textbf{constructive interference} 相长干涉).

\item exactly out of phase (crest meets \textbf{trough} 波谷, \textbf{phase difference} 相位差 $\pi$): the amplitude is $|A_{1} - A_{2}|$ (\textbf{destructive interference} 相消干涉).

\item any other phase difference $\phi$: the amplitude is somewhere between these two.

\end{itemize}

\par\medskip
\noindent\begin{minipage}{\linewidth}
\centering
\adjustbox{max width=\linewidth, max totalheight=0.42\textheight}{%
  \includegraphics{/Users/joschka/Desktop/Programming/eXercise/output/eXam/handouts/a_level_physics/assets/08-p150-constructive-interference.pdf}}\\[0.3em]
\emph{Two waves arriving in phase add to give double the amplitude (constructive)}
\end{minipage}\par\medskip

\par\medskip
\noindent\begin{minipage}{\linewidth}
\centering
\adjustbox{max width=\linewidth, max totalheight=0.42\textheight}{%
  \includegraphics{/Users/joschka/Desktop/Programming/eXercise/output/eXam/handouts/a_level_physics/assets/08-p151-destructive-interference.pdf}}\\[0.3em]
\emph{Two waves arriving exactly out of phase cancel to zero (destructive)}
\end{minipage}\par\medskip

For \textbf{intensity} 强度, $I \propto A^{2}$. Two equal waves meeting in phase give intensity $(2A)^{2} = 4A^{2}$ — \textbf{four times} the intensity of one wave alone.

\section*{Stationary (standing) waves}

When two identical \textbf{progressive waves} 行波 travel in \textbf{opposite directions} and overlap, they make a \textbf{stationary wave} 驻波. Examples: a wave on a string \textbf{reflected} 反射 from a fixed end overlapping the incoming wave; sound in an air column reflected from a closed end; microwaves between an emitter and a metal sheet.

\par\medskip
\noindent\begin{minipage}{\linewidth}
\centering
\adjustbox{max width=\linewidth, max totalheight=0.42\textheight}{%
  \includegraphics{/Users/joschka/Desktop/Programming/eXercise/output/eXam/handouts/a_level_physics/assets/08-p160-stationary-wave-formation.pdf}}\\[0.3em]
\emph{A stationary wave forms where two opposite waves overlap (N marks a node, A an antinode)}
\end{minipage}\par\medskip

\subsection*{Nodes and antinodes}

In a stationary wave:

\begin{itemize}
\item \textbf{node} 波节 — a point that is always at zero displacement (the two waves always cancel). The distance between next-door nodes is $\lambda/2$.

\item \textbf{antinode} 波腹 — a point of largest amplitude (the two waves always add). The distance between next-door antinodes is $\lambda/2$.

\item a node and the next antinode are $\lambda/4$ apart.

\end{itemize}

Particles between two nodes oscillate \textbf{in phase} with each other, but with different amplitudes (largest at the antinode, zero at the nodes). Particles on opposite sides of a node oscillate in \textbf{antiphase} 反相 (phase difference $\pi$).

\par\medskip
\noindent\begin{minipage}{\linewidth}
\centering
\adjustbox{max width=\linewidth, max totalheight=0.42\textheight}{%
  \includegraphics{/Users/joschka/Desktop/Programming/eXercise/output/eXam/handouts/a_level_physics/assets/08-p158-string-fundamental-mode.pdf}}\\[0.3em]
\emph{Fundamental mode on a stretched string — one loop, with L equal to half a wavelength}
\end{minipage}\par\medskip

How a stationary wave differs from a progressive wave: it does not carry \textbf{energy} 能量 along its length, the pattern does not move along, and the nodes stay fixed; a progressive wave has the same amplitude everywhere and carries energy.

\subsection*{Measuring wavelength from node spacing}

Drive a string with a vibrator at frequency $f$ until a stationary pattern appears. Measure the distance between two well-separated nodes and divide by the number of half-\textbf{wavelengths} 波长 between them. Then $\lambda$ is known, and $v = f\lambda$ gives the wave speed.

For a tube closed at one end and open at the other (a \textbf{resonance tube} 共鸣管), the closed end is a displacement node and the open end is a displacement antinode. The \textbf{fundamental} 基频 has $L = \lambda/4$; the next resonance is at $L = 3\lambda/4$; and so on. For a tube open at both ends, both ends are antinodes; the fundamental is at $L = \lambda/2$.

\par\medskip
\noindent\begin{minipage}{\linewidth}
\centering
\adjustbox{max width=\linewidth, max totalheight=0.42\textheight}{%
  \includegraphics{/Users/joschka/Desktop/Programming/eXercise/output/eXam/handouts/a_level_physics/assets/08-p161-closed-pipe-fundamental.pdf}}\\[0.3em]
\emph{Fundamental mode in a closed pipe — a node at the closed end, an antinode at the open end}
\end{minipage}\par\medskip

\section*{Diffraction}

\textbf{Diffraction} 衍射 is the spreading of a wave after it passes through a gap or around an \textbf{obstacle} 障碍物. All waves diffract — water, sound, light, microwaves.

The amount of spreading depends on the ratio of \textbf{wavelength to gap width}:

\begin{itemize}
\item gap \textbf{much wider} than $\lambda$: very little spreading; the wave goes nearly straight through.

\item gap \textbf{about the size of} $\lambda$: a lot of spreading; the wave fans out.

\item gap \textbf{smaller than} $\lambda$: very strong spreading; the gap acts almost like a point source.

\end{itemize}

Show this with water waves in a \textbf{ripple tank} 水波槽: straight waves meet a barrier with a gap, and the waves curve more as the gap is made narrower. The same idea is why you can hear someone around a corner (speech has $\lambda$ near 1 m, close to the gap size) but cannot see them (visible light has $\lambda \sim 500\ \text{nm}$, far smaller than the gap).

\par\medskip
\noindent\begin{minipage}{\linewidth}
\centering
\adjustbox{max width=\linewidth, max totalheight=0.42\textheight}{%
  \includegraphics{/Users/joschka/Desktop/Programming/eXercise/output/eXam/handouts/a_level_physics/assets/08-p165-ripple-tank-diffraction.pdf}}\\[0.3em]
\emph{Diffraction in a ripple tank — a wide gap (a) spreads the waves little, a narrow gap (b) much more}
\end{minipage}\par\medskip

\section*{Interference}

\textbf{Interference} 干涉 is the superposition of two \textbf{coherent} 相干 waves to give a steady pattern of high-amplitude regions (constructive) and low-amplitude regions (destructive).

\par\medskip
\noindent\begin{minipage}{\linewidth}
\centering
\adjustbox{max width=\linewidth, max totalheight=0.42\textheight}{%
  \includegraphics{/Users/joschka/Desktop/Programming/eXercise/output/eXam/handouts/a_level_physics/assets/08-p150-two-source-interference.pdf}}\\[0.3em]
\emph{Two coherent sources give lines of maximum displacement where crests meet crests}
\end{minipage}\par\medskip

\subsection*{Coherence}

Two sources are \textbf{coherent} when they emit waves with a \textbf{constant phase difference} (which also needs the same frequency). Two separate lamps are not coherent — their phase changes randomly, so any pattern flickers too fast to see and you get only an average.

To make coherent light from one source, pass it through two \textbf{slits} 狭缝 in a \textbf{double-slit} 双缝 setup. Both slits are lit by the same wavefront, so the two beams keep a fixed phase relationship.

\subsection*{Conditions for a clear pattern}

To see two-source fringes you need:

\begin{enumerate}
\item two \textbf{coherent} sources (constant phase difference).

\item roughly equal amplitudes (or the dark regions are not very dark).

\item the waves overlap where you look.

\item for light (a transverse wave), the same plane of \textbf{polarisation} 偏振.

\end{enumerate}

\subsection*{Path difference}

For two coherent sources, what happens at a point depends on the \textbf{path difference} 路程差 $\Delta x$ between the two waves arriving there:

\begin{itemize}
\item constructive: $\Delta x = n\lambda$ (for whole numbers $n = 0, 1, 2, \ldots$).

\item destructive: $\Delta x = (n + \tfrac{1}{2})\lambda$.

\end{itemize}

\subsection*{Double-slit (Young's) experiment}

For two slits a distance $a$ apart, with a screen a distance $D$ away (assume $D \gg a$), light of wavelength $\lambda$ makes fringes on the screen.

\par\medskip
\noindent\begin{minipage}{\linewidth}
\centering
\adjustbox{max width=\linewidth, max totalheight=0.42\textheight}{%
  \includegraphics{/Users/joschka/Desktop/Programming/eXercise/output/eXam/handouts/a_level_physics/assets/08-p154-youngs-double-slit.pdf}}\\[0.3em]
\emph{Young's double-slit experiment — the single slit makes the two slits coherent sources}
\end{minipage}\par\medskip

The \textbf{fringe spacing} 条纹间距 $x$ (one \textbf{fringe} 条纹 to the next) is

\[
\lambda = \frac{a x}{D}, \qquad x = \frac{\lambda D}{a}.
\]

Bright fringes (\textbf{maximum} 极大) are where the path difference is a whole number of $\lambda$; dark fringes (\textbf{minimum} 极小) where it is $(n + \tfrac{1}{2})\lambda$. The fringes are equally spaced.

To make the fringe spacing smaller: increase $a$ (slits further apart), reduce $D$ (screen closer), or use a shorter $\lambda$ (bluer light).

\section*{Diffraction grating}

A \textbf{diffraction grating} 衍射光栅 has many equally spaced slits — often hundreds or thousands per millimetre. Each slit is a coherent source. A maximum is seen at angle $\theta$ from the \textbf{normal} 法线 to the grating when

\[
d \sin\theta = n \lambda,
\]

where $d$ is the \textbf{slit spacing} 缝间距 (centre to centre), $n = 0, \pm 1, \pm 2, \ldots$ is the \textbf{order} 级次, and $\lambda$ is the wavelength.

Compared with the double slit, a grating gives much \textbf{sharper} maxima, because more slits add together — every other direction is cancelled by many slits.

\par\medskip
\noindent\begin{minipage}{\linewidth}
\centering
\adjustbox{max width=\linewidth, max totalheight=0.42\textheight}{%
  \includegraphics{/Users/joschka/Desktop/Programming/eXercise/output/eXam/handouts/a_level_physics/assets/08-p168-diffraction-grating.pdf}}\\[0.3em]
\emph{A diffraction grating splits monochromatic light into sharp maxima on a screen}
\end{minipage}\par\medskip

\subsection*{Slit spacing from "lines per mm"}

If a grating has $N$ lines per millimetre, then $d = 1/N$ millimetres $= 10^{-3}/N$ metres. For $450$ lines per mm, $d = 1/450\ \text{mm} \approx 2.22\ \text{µm}$.

\subsection*{Highest order}

For a given grating and wavelength, $\sin\theta = n\lambda/d$ cannot be more than $1$, so the highest order seen is

\[
n_{\text{max}} = \left\lfloor \frac{d}{\lambda} \right\rfloor.
\]

If $d/\lambda = 3.27$, orders up to $n = 3$ exist; $n = 4$ would need $\sin\theta > 1$ and is not seen.

\subsection*{Finding $\lambda$ with a grating}

Shine parallel light of unknown wavelength straight at the grating. Measure the angle $\theta_{1}$ of the first-order maximum from the centre. Then $\lambda = d \sin\theta_{1}$. Repeating for higher orders and averaging reduces error.


\end{document}
